% Ad M. Brutum 1.15 --- headnote
% ad-brutum-01-15, mid-early July 43 BC, Rome

\textit{Cicero to M. Brutus, from Rome, mid-early July
43 BC --- Perseus dateline \textit{Scr. Romae circ. med.
in. Quint. a. 711 (43)}, ``approximately the middle of
early Quintilis,'' i.e. roughly 5--14 July; Shackleton
Bailey places it about 14 July. By this point the
Mutina war has been won and then thrown away: the two
consuls, Hirtius and Pansa, are dead; D. Brutus is
isolated in Cisalpine Gaul; Lepidus has joined his
forces to Antony's in Narbonese Gaul (the ``crime of
Lepidus'' of 1.14); and Octavian, who turns twenty
this autumn, is the only field commander still under
the Senate's nominal control --- and increasingly not.
Brutus has written a letter critical of Cicero's
policy, and Messala (the future patron of Tibullus and
Ovid, here making his first major appearance in
Cicero's correspondence) is carrying this answer
eastward to him.}

\textit{The letter is a sustained political
self-defence built on a frame Cicero attributes to
Solon: \textit{rem publicam contineri duabus rebus,
praemio et poena} --- the commonwealth is held together
by two things, reward and punishment. The text of the
ascription is corrupt (the daggered crux at
\textit{neque solum ut Solonis dictum usurpem}); the
maxim itself is not in fact in Solon and may be
Cicero's own. Brutus has accused him of being prodigal
in honours; Cicero defends each grant in turn --- to
the young Caesar, to D. Brutus, to L. Plancus, to the
dead consuls Hirtius and Pansa and to Aquila. The
single concession is the statue of Lepidus on the
Rostra (\textit{nos illum honore studuimus a furore
revocare. vicit amentia levissimi hominis nostram
prudentiam}). On punishment Cicero is uncompromising:
this war is unlike any earlier civil war --- for the
losers there will be no commonwealth at all. The
letter closes with the same demand as 1.14, now
sharpened: \textit{sed propera, per deos! scis quantum
sit in temporibus, quantum in celeritate.} Brutus
will not come. This is the third-from-last surviving
letter of the entire \textit{Ad Brutum} correspondence.}

\textit{The architecture is a Ciceronian set-piece
turned to political use: the rewards-and-punishments
frame announced in section 3, the ``honours'' catalogue
(sections 7--9), the ``punishment'' catalogue (sections
10--11), and the closing plea (section 12). The
Themistocles citation in section 11 --- ``even the
children of Themistocles were in want'' --- defends
inherited proscription against a sentimental objection
by reaching for the standing example from Athenian
history.}
